\documentclass[11pt]{article}
\usepackage[spanish]{babel}
\usepackage[utf8]{inputenc}
\usepackage[T1]{fontenc}
\usepackage{amsmath}
\usepackage{booktabs}
\usepackage{graphicx}
\usepackage{hyperref}
\usepackage[backend=biber,style=numeric]{biblatex}
\addbibresource{bibliografia.bib}

\title{Analizador de Notas Musicales mediante FFT \\ \large Proyecto Integrador 6 --- Fase 1}
\author{Monica Daniela Alvarenga Mejia \and Rodrigo José Marroquín Quijada}
\date{Instituto Kriete de Ingeniería y Ciencias --- Señales y Sistemas (EE2004)}

\begin{document}
\maketitle

\section{Introducción}
% TODO: motivación, alcance de la Fase 1 y resumen de resultados (depende de
% los resultados finales; redactar al cierre).

\section{Modelo matemático}
\label{sec:modelo}

\subsection{Muestreo y teorema de Nyquist}
La señal continua $x(t)$ captada por el micrófono se muestrea a
$f_s = 48\,000$~Hz, produciendo la secuencia discreta $x[n] = x(nT_s)$, con
$T_s = 1/f_s$. Las doce notas de la octava Do4--Si4 tienen frecuencias
fundamentales entre 261{,}63~Hz y 493{,}88~Hz. Considerando los armónicos de
interés hasta unos 4~kHz, el teorema de muestreo de Nyquist--Shannon exige
$f_s > 2 \cdot 4\,000 = 8\,000$~Hz; los 48~kHz elegidos dan un margen de más
de cinco veces ese contenido espectral, lo que permite además observar el
caso --- frecuente en instrumentos reales --- en que un armónico aparece con
mayor amplitud que la fundamental \cite{oppenheim2010,rossing1990science}.

\subsection{Transformada discreta de Fourier (DFT/FFT)}
Sobre cada segmento preprocesado de $N$ muestras se calcula la DFT

\begin{equation}
  X[k] = \sum_{n=0}^{N-1} x[n]\, e^{-j 2\pi k n / N}, \qquad k = 0, \dots, N-1,
\end{equation}

evaluada de forma eficiente mediante el algoritmo FFT (\texttt{np.fft.rfft},
que aprovecha la simetría hermítica de una señal real). El eje de
frecuencias correspondiente a cada bin $k$ es

\begin{equation}
  f_k = \frac{k \cdot f_s}{N},
\end{equation}

lo que con $N = 32\,768$ da una resolución de bin
$\Delta f = f_s/N \approx 1{,}465$~Hz \cite{oppenheim2010,smith2007mdft}.

\subsection{Ventaneo de Hann}
Analizar un segmento finito de la señal equivale a multiplicarlo por una
ventana rectangular, lo que produce fuga espectral (\emph{spectral leakage})
y dificulta separar la fundamental de picos vecinos. Se aplica en su lugar
una ventana de Hann,

\begin{equation}
  w[n] = 0{,}5 \left(1 - \cos\!\left(\frac{2\pi n}{N-1}\right)\right),
\end{equation}

que reduce los lóbulos laterales a costa de ensanchar ligeramente el lóbulo
principal, un compromiso favorable para localizar picos aislados
\cite{harris1978windows}.

\subsection{Interpolación parabólica del pico}
% TODO (17/09): una vez implementada en preprocess/fft_pipeline, confirmar
% que la fórmula usada coincide exactamente con esta y referenciar el
% módulo correspondiente.
La resolución de bin $\Delta f \approx 1{,}465$~Hz no es suficiente por sí
sola para cumplir el criterio de error de 5~cents en Do4 (que equivale a
$\approx 0{,}76$~Hz). Se refina la ubicación del pico ajustando una parábola
a la magnitud (en dB) del bin máximo $\beta$ y sus vecinos $\alpha$
(izquierdo) y $\gamma$ (derecho):

\begin{equation}
  \delta = \frac{1}{2} \cdot \frac{\alpha - \gamma}{\alpha - 2\beta + \gamma},
  \qquad
  f_0 = (k + \delta) \cdot \frac{f_s}{N},
\end{equation}

donde $k$ es el índice del bin de máxima magnitud \cite{smith2007mdft}.

\subsection{Identificación de la nota y error en cents}
La frecuencia de referencia de cada nota $n$ semitonos por encima o debajo
de La4 sigue el temperamento igual,

\begin{equation}
  f_n = 440 \cdot 2^{(n-49)/12} \; \text{Hz}, \qquad \text{La4} = 440~\text{Hz}
  \;\; \cite{iso16-1975},
\end{equation}

y el error de la frecuencia detectada $f_0$ respecto de la tabla se expresa
en cents (1/100 de semitono):

\begin{equation}
  \text{error}_{\text{cents}} = 1200 \cdot \log_2\!\left(\frac{f_0}{f_{\text{tabla}}}\right).
\end{equation}

\section{Metodología}
\label{sec:metodologia}

\subsection{Adquisición}
Cada una de las doce notas de la octava Do4--Si4 se graba de forma
individual con el micrófono y la tarjeta de sonido de la laptop, en
archivos WAV de aproximadamente 1~s de duración. Los primeros 50--100~ms de
cada grabación (el ataque del martillo sobre la cuerda, con transitorios de
alta energía y contenido espectral no estacionario) se descartan antes del
análisis. Las grabaciones originales se conservan sin modificar en
\texttt{audio/originales/} como evidencia de adquisición; todo el
procesamiento se realiza sobre copias en \texttt{audio/procesados/}.

\subsection{Pipeline de procesamiento}
El pipeline es el mismo para los tonos sintéticos (usados para validar el
código antes de introducir señales reales) y para las grabaciones reales:

\begin{enumerate}
  \item \textbf{Preprocesamiento}: conversión a mono, remoción de offset DC,
    recorte del ataque, ventana de Hann (Sección~\ref{sec:modelo}).
  \item \textbf{FFT}: \texttt{np.fft.rfft} sobre la señal ventaneada; eje de
    frecuencias con \texttt{np.fft.rfftfreq(N, 1/fs)}.
  \item \textbf{Detección de pico}: máximo de magnitud en la banda
    80--1000~Hz, que cubre las fundamentales de Do4 a Si4 sin confundirlas
    con armónicos graves de notas más bajas.
  \item \textbf{Refinamiento sub-bin}: interpolación parabólica
    (Sección~\ref{sec:modelo}).
  \item \textbf{Identificación y error}: comparación contra la tabla de
    temperamento igual y cálculo del error en cents.
\end{enumerate}

Un tono sintético con frecuencia conocida sirve como caso de control: si el
pipeline falla sobre él, el error está en el código; si el sintético pasa y
la grabación real falla, el problema es de captura o de la señal física, no
del algoritmo.

\subsection{Validación}
% TODO: completar con los resultados reales (sintéticos con interpolación +
% grabaciones reales) cuando estén disponibles. No adelantar números aquí.

\section{Resultados}
% TODO: pendiente de la interpolación parabólica y de las grabaciones
% reales verificadas en Audacity.

\section{Discusión}
% TODO.

\printbibliography

\end{document}
